В результате лабораторной работы было вычислено значение определенного интеграла с помощью различных численных методов: метода прямоугольников, метода трапеций, метода Симпсона, также использовался метод Ричардсона для уточнения результата. Все методы были применены с целью достижения заданной погрешности $\epsilon = 10^{-3}$.

Метод Симпсона оказался наиболее эффективным, требуя наименьшее количество разбиений для достижения заданной погрешности. Методы прямоугольников и трапеций показали одинаковые результаты в данном случае.